This thesis investigated whether a decision-support system combining
ontology-based knowledge representation, semantic retrieval from
scientific literature, and structured implementation support can
provide relevant and justified machine learning method recommendations
for users with limited machine learning expertise,
and support the first steps of implementing the chosen method. 
The question was
examined through the design and evaluation of MLguide, a web-based
system built within a Design Science Research framework. 
MLguide draws on an ontology of machine learning knowledge and a database of 
research articles. For output it produces a ranked list of machine learning 
methods, a set of supporting articles, and a starter notebook. The system was evaluated
through user testing in two machine learning course projects and through an
example-based evaluation on constructed problems.

The evaluation showed that the recommended methods matched the
submitted problems well across different ML areas, and that the three
elements function together in a working system. Several properties in
the ontology have low coverage, so it contributes little to the
ranking of methods, which is instead driven by the semantic retrieval. The
recommendations are grounded in the literature rather than in fixed
rules, and the articles shown for each method make the basis for a
recommendation visible to the user. The weaker point is justification,
since a method counts as supported whenever a matching article
mentions it, even when the method is not the main contribution of the
article. The implementation support is real but limited, since the starter
notebooks are simple and cover only part of the method set, though
the structure makes them easy to extend.

The main contribution of this thesis is MLguide itself, a working
system that brings together three capabilities that earlier work has
kept apart. MLguide combines structured method
knowledge, literature-based recommendation, and a starting point for
implementation in a single system, and presents the articles behind
each recommendation so that the user can judge it. A further
contribution is the link between the ontology and the articles, where
each article is classified and connected to the methods it mentions as
part of the retrieval pipeline. This lets the method knowledge grow as
the article base grows, and keeps the system open to a different domain by
replacing the article base. In Design Science Research terms, the
artifact also yields design knowledge that can inform similar systems,
in particular the value of pairing structured method knowledge with
retrieval from the literature, and the trade-offs this involves.

The clearest path forward builds on the
link between the articles and the ontology that already drives the
recommendations. Extending this link to draw more structured knowledge from the
literature automatically would build up the knowledge in the ontology
while reducing the manual effort behind it.
The starter notebooks could also be developed further, both to cover more
methods, including reinforcement learning, and to let the user configure
the method directly, with suggested parameter values.
The system would also benefit from a broader evaluation, across several
machine learning courses and towards engineers in industry, which would
give firmer conclusions than were possible in this thesis.